% chapters/testing/unit_tests.tex
% Sección 6.2: Pruebas unitarias (~1 página)
% ============================================================================
%
% GUÍA: Describir las pruebas unitarias realizadas.
%       Referencia a \ref{rnf:test:unit}.
%
%   1. Qué se prueba unitariamente:
%      - Máquina de estados del examen (transiciones válidas e inválidas).
%      - Evaluación de reglas de asignación de correctores.
%      - Cálculo de notas (agregación de anotaciones).
%      - Verificación de permisos (cadena de evaluación).
%      - Identificación por similitud (algoritmo OCR).
%      - Validaciones de serializers.
%
%   2. Ejemplo representativo (incluir un test como código):
%      \PythonCode[cod:unit_test]{Test unitario de ejemplo}
%                 {Prueba unitaria de la transición de estado del examen}
%                 {codes/test_exam_state.py}{1}{20}{1}
%      TODO: Copiar un test representativo a codes/test_exam_state.py.
%
%   3. Resultados resumidos:
%      - Número de tests unitarios.
%      - Porcentaje de cobertura de la lógica de negocio.
%      - Tabla resumen por módulo (si hay suficientes datos):
%
%      \begin{table}{tab:unit_results}{Resumen de pruebas unitarias por módulo}
%        \begin{tabular}{l r r}
%          \textbf{Módulo} & \textbf{Tests} & \textbf{Cobertura (\%)} \\
%          \hline
%          Autenticación      & --  & -- \\
%          Permisos           & --  & -- \\
%          Máquina de estados & --  & -- \\
%          OCR / similitud    & --  & -- \\
%          ...                & --  & -- \\
%          \hline
%          \textbf{Total}     & --  & -- \\
%        \end{tabular}
%      \end{table}
%
%   NOTA: Resultados detallados (logs, listado completo de tests) van en
%   el \cref{ap:detailed_results}.
% ============================================================================

% TODO: Redactar la sección de pruebas unitarias.
